% Generated from locale/tr/content/lambda-calculus/church-rosser/definitions-and-properties.tex; do not hand-edit.


\olfileid[tr]{lam}{cr}{dap}

\olsection{Tanım ve Özellikler}

Bu bölümde Church--Rosser özelliğini ve bu özelliğe ilişkin bazı yaygın
sonuçları tanıtıyoruz.

\begin{defn}[Church--Rosser özelliği, CR] Terimler üzerindeki bir
  $\xredone$ bağıntısı, $M\xredone P$ ve $M\xredone Q$ olduğunda
  $P\xredone N$ ve $Q\xredone N$ olacak biçimde bir $N$ her zaman varsa
  \emph{Church--Rosser özelliğini} sağlar.
\end{defn}

Lambda hesabını, normal biçimdeki terimlerin ``değerler'', terimler
üzerindeki indirgenebilirlik bağıntısının ise ``hesap kuralları'' olduğu bir
hesaplama modeli olarak görebiliriz. Church--Rosser özelliği, bir hesaplamayı
sürdürmenin birden çok yolu bulunduğunda bile ifadenin yalnızca tek bir son
değeri olduğunu söyler.

Temel cebirden bir örnek alalım. $4\times(1+2)+3$ ifadesi birden çok yoldan
hesaplanabilir: önce $1+2$ ifadesini $3$'e indirgersek $4\times3+3$;
dağılma özelliğini önce $4\times(1+2)$ ifadesine uygularsak
$4\times1+4\times2+3$ elde edilir. Bununla birlikte her ikisi de daha sonra
$12+3$ ifadesine indirgenebilir.

$\xredone$ bağıntısını $\beta$-indirgeme olarak alırsak Church--Rosser
özelliğinin şu sonucunu kolayca görürüz: bir terimin normal biçimi varsa bu
biçim tektir. Gerçekten $M$'nin, ikisi de normal biçimde olan $P$ ve $Q$'ya
indirgenebildiğini varsayalım. Church--Rosser özelliğine göre hem $P$ hem de
$Q$'nun indirgenebildiği bir $N$ vardır. $P$ ve $Q$ normal biçimde olduğundan
bunların $N$'ye indirgenmesi ancak önemsiz indirgeme olabilir; yani $P$, $Q$
ve $N$ özdeştir. Bu sonuç bir terimin \emph{normal biçiminden} tekil olarak
söz etmemizi haklı çıkarır.

Lambda hesabını bir hesaplama modeli olarak gördüğümüzde bir terimin normal
biçimi, o terimle başlayan hesaplamanın ``nihai sonucu'' sayılabilir.
Yukarıdaki sonuç, $4\times(1+2)+3$ hesabının tek sonucu olan $15$ gibi, bir
hesaplamanın varsa yalnızca tek bir nihai sonucu bulunduğu anlamına gelir.

\begin{thm} \ollabel{thm:str}
  Bir $\xredone$ bağıntısı Church--Rosser özelliğini sağlıyor ve $\xred$,
  $\xredone$ bağıntısını içeren en küçük geçişli bağıntıysa $\xred$ de
  Church--Rosser özelliğini sağlar.
\end{thm}

\begin{proof}
  Şunları varsayalım:
  \begin{align*}
    M & \xredone P_1 \xredone \dots \xredone P_m \text{ ve}\\
    M & \xredone Q_1 \xredone \dots \xredone Q_n.
  \end{align*}
  Teoremi yüksekliği $m+1$, genişliği $n+1$ olan bir $N$ terimler ızgarası
  kurarak ispatlayacağız. $i$'inci satır ile $j$'inci sütundaki terimi
  $N_{i,j}$ ile gösterelim.

  $N$ ızgarasını $N_{i,j}\xredone N_{i+1,j}$ ve
  $N_{i,j}\xredone N_{i,j+1}$ olacak biçimde kurarız. Tanım şöyledir:
  \begin{align*}
    N_{0,0} &= M \\
    N_{i,0} &= P_i && \text{$1\le i\le m$ ise} \\
    N_{0,j} &= Q_j && \text{$1\le j\le n$ ise} \\
  \intertext{ve diğer durumlarda:}
    N_{i,j} &= R.
  \end{align*}
  Burada $R$, $N_{i-1,j}\xredone R$ ve $N_{i,j-1}\xredone R$ olacak
  biçimde bir terimdir. $\xredone$ Church--Rosser özelliğini sağladığından
  böyle bir terim her zaman vardır.

  Şimdi $N_{m,0}\xredone\dots\xredone N_{m,n}$ ve
  $N_{0,n}\xredone\dots\xredone N_{m,n}$ elde edilir. $N_{m,0}$ teriminin
  $P$, $N_{0,n}$ teriminin de $Q$ olduğuna dikkat edin. Sonuç $\xred$
  bağıntısının tanımından çıkar.
\end{proof}

